% =====================================================================
%  CARNET D'ENTRAÎNEMENT — FICHE 8 — THÈME 1 — NIVEAU DIFFICILE — CORRIGÉ
% =====================================================================
\documentclass[11pt,a4paper]{article}
\usepackage[utf8]{inputenc}
\usepackage[T1]{fontenc}
\usepackage{lmodern}
\usepackage[french]{babel}
\usepackage{amsmath,amssymb,mathtools}
\usepackage{icomma}
\usepackage{siunitx}
\sisetup{output-decimal-marker={,},per-mode=power,group-separator={\,},group-minimum-digits=5}
\usepackage[version=4]{mhchem}
\usepackage{xcolor}
\usepackage{tikz}
\usepackage[most]{tcolorbox}
\usepackage{enumitem}
\usepackage{array}
\usepackage{fancyhdr}
\usepackage[hidelinks]{hyperref}
\usetikzlibrary{arrows.meta,positioning,calc}
\usepackage[a4paper,margin=2.1cm,headheight=26pt,headsep=10pt]{geometry}
\usepackage{titlesec}

\definecolor{primary}{RGB}{150,98,162}
\definecolor{primarydark}{RGB}{92,52,110}
\definecolor{excol}{RGB}{0,135,90}
\definecolor{exocol}{RGB}{200,60,40}

\tcbset{boxbase/.style={enhanced,breakable,boxrule=0.9pt,arc=2.5pt,left=3mm,right=3mm,top=2.3mm,bottom=2.3mm,fonttitle=\bfseries,coltitle=white,toptitle=0.6mm,bottomtitle=0.6mm}}
\newtcolorbox[auto counter]{correction}[1][]{boxbase,colback=excol!4,colframe=excol,title={\textsc{Corrigé}~\thetcbcounter\ \ifx\relax#1\relax\relax\else\textmd{--- \itshape #1}\fi}}

\newcommand{\dH}{\Delta H}
\newcommand{\dU}{\Delta U}
\newcommand{\rep}[1]{\textcolor{primarydark}{\textbf{#1}}}

\pagestyle{fancy}
\fancyhf{}
\fancyhead[L]{\small\textcolor{primarydark}{\textbf{Fiche 8} — Thème 1 : Premier principe et enthalpie}}
\fancyhead[R]{\small\textcolor{primarydark}{\textsc{Difficile — corrigé}}}
\fancyfoot[C]{\small\thepage}
\renewcommand{\headrulewidth}{0.8pt}
\renewcommand{\headrule}{\hbox to\headwidth{\color{primary}\leaders\hrule height \headrulewidth\hfill}}
\setlength{\parindent}{0pt}\setlength{\parskip}{3pt}

\begin{document}

\begin{center}
\begin{tikzpicture}
\node[rectangle,rounded corners=7pt,fill=excol,text=white,minimum width=16cm,
      minimum height=1.1cm,align=center,inner sep=6pt]
  {{\large\bfseries Thème 1 — Premier principe et enthalpie}\\[1pt]{\small Niveau \textsc{Difficile} — corrigé détaillé}};
\end{tikzpicture}
\end{center}

\begin{correction}[détente d'un gaz dans un piston — bilan complet]
\begin{enumerate}[label=\alph*),leftmargin=2em,topsep=2pt,itemsep=3pt]
  \item $\Delta V = (3{,}0-1{,}0)=\SI{2.0}{\litre}=\SI{2.0e-3}{\cubic\metre}$.
        \[ W = -p\,\Delta V = -(\SI{1.0e5}{\pascal})(\SI{2.0e-3}{\cubic\metre}) = \rep{\SI{-200}{\joule}}. \]
        $W<0$ : en se dilatant, le gaz \rep{cède} du travail (il pousse le piston).
  \item $\dU = Q + W = (+250) + (-200) = \rep{\SI{+50}{\joule}}$.
  \item $\dU>0$ : l'énergie interne a \rep{augmenté} de \SI{50}{\joule}. Interprétation : sur les
        \SI{250}{\joule} de chaleur reçue, \SI{200}{\joule} sont \emph{restitués} au milieu sous forme de
        travail (le piston monte), et seuls \SI{50}{\joule} restent stockés dans le gaz. \emph{La chaleur
        reçue ne sert donc pas entièrement à augmenter $U$.}
\end{enumerate}
\end{correction}

\begin{correction}[combustion du carbone — énergie et quantité de matière]
\begin{enumerate}[label=\alph*),leftmargin=2em,topsep=2pt,itemsep=2pt]
  \item $\dH<0$ : les \rep{réactifs} ($\ce{C}+\ce{O2}$) sont \emph{plus haut} que le produit ($\ce{CO2}$).
        La réaction est \rep{exothermique}.
  \item Énergie dégagée $= n\,|\dH| = 3{,}0 \times 394 = \rep{\SI{1182}{\kilo\joule}}$.
  \item $n = \dfrac{1970}{394} = \SI{5.0}{\mol}$, \; d'où $m = n\,M = 5{,}0 \times 12 = \rep{\SI{60}{\gram}}$
        de carbone.
  \item Par mole de carbone, la combustion \emph{complète} dégage \SI{394}{\kilo\joule}, contre seulement
        \SI{111}{\kilo\joule} pour la combustion \emph{incomplète}. La \rep{combustion complète} dégage
        beaucoup plus d'énergie : brûler jusqu'au \ce{CO2} exploite toute l'énergie du carbone, alors que
        s'arrêter au \ce{CO} en « gaspille » une grande partie (et produit un gaz toxique).
\end{enumerate}
\end{correction}

\begin{correction}[distinguer $\dU$ et $\dH$ — l'ordre de grandeur du travail]
\begin{enumerate}[label=\alph*),leftmargin=2em,topsep=2pt,itemsep=2pt]
  \item $\Delta V = \SI{+2.0}{\litre} = \SI{2.0e-3}{\cubic\metre}$, d'où
        \[ W = -p\,\Delta V = -(\SI{1.0e5}{\pascal})(\SI{2.0e-3}{\cubic\metre}) = \SI{-200}{\joule}
             = \rep{\SI{-0.2}{\kilo\joule}}. \]
  \item $\dU = \dH + W = -180 + (-0{,}2) = \rep{\SI{-180.2}{\kilo\joule}}$.
  \item $\dU \approx \dH$ : le travail (\SI{0.2}{\kilo\joule}) est \rep{négligeable} devant la chaleur de
        réaction (\SI{180}{\kilo\joule}), soit environ $0{,}1\,\%$. \emph{À retenir :} pour une réaction
        chimique, $\dU$ et $\dH$ sont très proches ; c'est pourquoi la thermochimie raisonne
        commodément sur $\dH$ (la chaleur à pression constante).
\end{enumerate}
\end{correction}

\end{document}
